\originalchapter{本科作业: 介质、磁场与守恒}
本章解答均为\textbf{AI 独立解答}. 所列缺失项在官方作业 sheet 中存在, 但只给教材题号或不完整的追加子问; 原 sheet 已归档不等于已取得商业教材的完整题面.

\originalsection{作业6}
\begin{exercise}
原题4 (15分). 两个理想电偶极, $\pp_1=p_1\hat{\vct z}$ 位于原点, $\pp_2=p_2\hat{\vct x}$ 位于 $r\hat{\vct x}$. (a) 求彼此作用力, 是否等大反向且沿连心线? (b) 加入各自的 $\pp\times\ee$ 力矩, 核对系统总力矩.
\end{exercise}
\begin{solution}
两偶极相互作用能为
$U=k[\pp_1\cdot\pp_2-3(\pp_1\cdot\nn)(\pp_2\cdot\nn)]/r^3$, $k=1/(4\pi\eps)$. 为算力, 不能先把 $\pp_1\cdot\nn=0$ 代成处处恒零: 观察位置偏离 $x$ 轴时它会改变. 在 $(r,0,z)$ 附近 $U=-3kp_1p_2z/r^4+O(z^2)$, 故 $\vct F_2=3kp_1p_2\hat{\vct z}/r^4$, $\vct F_1=-\vct F_2$. 它们不沿连心线. 第2偶极处 $\ee_1=-kp_1\hat{\vct z}/r^3$, 第1处 $\ee_2=2kp_2\hat{\vct x}/r^3$, 因而 $\vct\tau_2=kp_1p_2\hat{\vct y}/r^3$, $\vct\tau_1=2kp_1p_2\hat{\vct y}/r^3$. 关于原点的轨道力矩为 $r\hat{\vct x}\times\vct F_2=-3kp_1p_2\hat{\vct y}/r^3$, 三项之和为零.
\end{solution}
\begin{exercise}
原题6 (15分). 两个半径 $R$ 的圆形极板位于 $z=\pm d/2$, 电压差 $V_0$ 固定. 厚度近似 $d$、边长大于 $2R$ 的方介质板部分插入, 左边界为 $x=x_0$, 介质占据 $x\ge x_0$. 求力随 $x_0$ 的变化.
\end{exercise}
\begin{solution}
作为附加近似模型, 取线性各向同性介质、$d\ll R$, 板厚近似充满间隙, 忽略边缘场. 介电常数 $\varepsilon>\eps$ 是这里明确补出的普通介质假设, 原题未给其数值. 圆盘中被填充面积为
$A_d=R^2\arccos(x_0/R)-x_0\sqrt{R^2-x_0^2}$, $|x_0|<R$; 这是积分 $\int_{x_0}^R2\sqrt{R^2-x^2}\dd x$. 导数为 $-2\sqrt{R^2-x_0^2}$. 固定电压的力为
$F_x=(V_0^2/2)C'(x_0)=-(\varepsilon-\eps)V_0^2\sqrt{R^2-x_0^2}/d$, 方向使 $x_0$ 减小、更多介质进入. 完全未插入或完全插入时, 此近似给 $F_x=0$. 若实际介质厚度为 $h<d$, 填充区单位面积电容应为 $[h/\varepsilon+(d-h)/\eps]^{-1}$, 把上式的 $(\varepsilon-\eps)/d$ 换成该量减 $\eps/d$ 即保留真空间隙效应. 极板间距不是小量时须另外求边缘场, 这里不把能量近似当成精确几何解.
\end{solution}
原题1、2、3、5分别只列 Griffiths 4.2、4.16、4.32、Jackson 4.10. 第1题虽有刚性电子云提示, 仍无原题完整参数, 不补造.

\originalsection{作业7}
\begin{exercise}
原题3 (15分). 半径 $R$ 的密绕螺线管从 $z_1$ 延伸到 $z_2>z_1$, 每单位长度 $n$ 匝, 电流从 $+z$ 看逆时针. (a) 求原点处磁场; (b) 求 $z_1\to-\infty,z_2\to\infty$ 极限.
\end{exercise}
\begin{solution}
每匝对原点贡献 $\mu_0IR^2\hat{\vct z}/[2(R^2+z'^2)^{3/2}]$. 乘 $n\dd z'$ 积分, 得
\[
\bb(0)=\frac{\mu_0nI}2\left[\frac{z_2}{\sqrt{R^2+z_2^2}}-
\frac{z_1}{\sqrt{R^2+z_1^2}}\right]\hat{\vct z}.
\]
极限括号为2, 因而无限长管内 $\bb=\mu_0nI\hat{\vct z}$. 本式对原点在管内外都适用; 管外两端相消使场较小.
\end{solution}
\begin{exercise}
原题7 (20分及10分加分). 验证 $\partial_i(x_j/r^3)=\PV(\delta_{ij}-3n_in_j)/r^3+(4\pi/3)\delta_{ij}\delta^3$. (a) 推出 $\nabla^2(1/r)$; (b) 推出精确电偶极场; (c) 推出精确磁偶极场; (d) 用小球积分检查接触项; (e) 对任意测试函数证明分布恒等式.
\end{exercise}
\begin{solution}
(e) 的完整证明见第4章: 对挖洞区域分部积分, 小球面项为 $\int n_in_j\varphi(a\nn)\dd\Omega\to(4\pi/3)\delta_{ij}\varphi(0)$. 剩下的是球形主值积分, 常数项角平均为零. (d) 将测试函数在原点附近取常数, 同一球面积分核对系数. (a) 缩并 $i=j$ 后主值项的迹为零, 左边为 $\nabla\cdot(\rr/r^3)=-\nabla^2(1/r)$, 接触项为 $4\pi\delta^3$. (b) 对 $V=p_jx_j/(4\pi\eps r^3)$ 取负梯度, 得第4章电偶极式. (c) $A_k=\mu_0\epsilon_{k\ell m}m_\ell x_m/(4\pi r^3)$, 取旋度后
\[
B_i=\frac{\mu_0}{4\pi}\left[m_i\partial_j(x_j/r^3)-m_j\partial_j(x_i/r^3)\right].
\]
第一项给 $4\pi m_i\delta^3$, 第二项给主值与 $4\pi m_i\delta^3/3$, 相减产生 $2\mu_0\mm\delta^3/3$.
\end{solution}
原题1、2、4、6只列 Griffiths 5.3、5.8、5.24、5.58. 第5题(a)--(c)仅引用5.43, MIT 追加(d)虽可读但依赖未给出的守恒量 $\vct Q$, 本稿列为不完整题面, 不把重新设定的单极模型冒充该题.

\originalsection{作业8}
\begin{exercise}
原题5 (20分), MIT sheet 注明基于 Griffiths 6.9 与 Jackson 5.19. 圆柱磁体半径 $a$, 长 $L$, 占据 $-L/2<z<L/2$, 均匀永磁化 $\vct M=M_0\hat{\vct z}$. (a) 求轴线各点的 $\vct H,\bb$; (b) 给出 $L/a=5$ 的 $|B|/(\mu_0M_0)$ 与 $|H|/M_0$ 曲线; (c) 求束缚电流; (d) 求远场偶极矩.
\end{exercise}
\begin{solution}
(c) 体电流为零, 侧面 $\vct K_b=M_0\hat{\vct\phi}$, 端面无面电流. 它等价于每单位长度 $nI=M_0$ 的螺线管. (a) 令
\[
f(z)=\frac12\left[\frac{z+L/2}{\sqrt{a^2+(z+L/2)^2}}
-\frac{z-L/2}{\sqrt{a^2+(z-L/2)^2}}\right],
\]
则 $\bb=\mu_0M_0f(z)\hat{\vct z}$; 管内 $\vct H=M_0[f(z)-1]\hat{\vct z}$, 管外 $\vct H=M_0f(z)\hat{\vct z}$. $\bb$ 的法向连续而 $\vct H$ 的法向因磁化表面变化而跳跃. (d) $\mm=\pi a^2LM_0\hat{\vct z}$. (b) 下图按上述公式采样, 不用有限采样证明场方程.
\begin{center}
\begin{tikzpicture}
\begin{axis}[width=.78\textwidth,height=5cm,xlabel={$z/a$},ylabel={归一化场强},xmin=-6,xmax=6,ymin=0,ymax=1,legend style={at={(.5,1.02)},anchor=south,legend columns=2},grid=major]
\addplot[thick,blue,domain=-6:6,samples=141] {.5*((x+2.5)/sqrt(1+(x+2.5)^2)-(x-2.5)/sqrt(1+(x-2.5)^2))};
\addlegendentry{$|B|/(\mu_0M_0)$}
\addplot[thick,dashed,red,domain=-2.5:2.5,samples=61] {1-.5*((x+2.5)/sqrt(1+(x+2.5)^2)-(x-2.5)/sqrt(1+(x-2.5)^2))};
\addlegendentry{$|H|/M_0$}
\addplot[thick,dashed,red,domain=-6:-2.5,samples=45] {.5*((x+2.5)/sqrt(1+(x+2.5)^2)-(x-2.5)/sqrt(1+(x-2.5)^2))};
\addplot[thick,dashed,red,domain=2.5:6,samples=45] {.5*((x+2.5)/sqrt(1+(x+2.5)^2)-(x-2.5)/sqrt(1+(x-2.5)^2))};
\end{axis}
\end{tikzpicture}\par
$L/a=5$, $M_0>0$. 在两端面取单侧极限, 不以曲线连接跳跃值.
\end{center}
\end{solution}
原题1--4、6、7只列 Griffiths 5.58、5.37、5.49、6.5、6.17、6.25. 第1题是作业7题6的延后布置, 不重复记为新增完整题面.

\originalsection{作业9}
\begin{exercise}
原题8 (20分加分). 记 $\Theta$ 为 Heaviside 阶跃函数. (a) 证明 $\Theta'=\delta$; (b) 证明正整数 $n$ 下 $\Theta^n=\Theta$; (c) 指出形式链式法则与分布乘积的矛盾, 求 $(\Theta^n)'$; (d) 证明 $f\delta_{x_0}=f(x_0)\delta_{x_0}$; (e) 对(d)微分并核对两边; (f) 对 $a>0$ 求 $\int\varphi(x)\delta(x^2-a^2)\dd x$.
\end{exercise}
\begin{solution}
(a) $\langle\Theta',\varphi\rangle=-\int_0^\infty\varphi'\dd x=\varphi(0)$. (b) 两函数只可能在原点取值不同, 单点不影响普通积分, 因而定义同一个分布. (c) 若误用链式法则, $n=2$ 给 $2\Theta\delta=\delta$, $n=3$ 给 $3\Theta^2\delta=\delta$; 再误用 $\Theta^2=\Theta$ 相乘, 会同时要求 $\Theta\delta=\delta/2$ 和 $\delta/3$. 原因是该乘积未按通常分布理论定义. 正确地先用(b)再求导, 得 $(\Theta^n)'=\delta$. (d) 与测试函数相乘取值即可. (e) 光滑函数乘分布合法, 且 $f\delta'=f(x_0)\delta'-f'(x_0)\delta$; 其中 $\langle f\delta',\varphi\rangle=-f'(x_0)\varphi(x_0)-f(x_0)\varphi'(x_0)$. 原(d)微分后的左边 $f'\delta+f\delta'$ 中第一项抵消 $-f'(x_0)\varphi(x_0)$, 故整个导数等式两边皆为 $-f(x_0)\varphi'(x_0)$. (f) 两个简单根的导数绝对值皆为 $2a$, 故 $[\varphi(a)+\varphi(-a)]/(2a)$. $a=0$ 是重根, 不能直接代入简单根公式.
\end{solution}
原题1--7只列 Griffiths 7.8、7.20、7.42、7.49、7.58、7.59、8.1. 本稿第7、8、11章补充相应概念, 但不把这些补充例子标为已取得的教材原题解答.

\originalsection{作业10}
\begin{exercise}
原题3 (25分). 半径 $R$ 的均匀带电球壳带电 $Q$, 以 $\omega\hat{\vct z}$ 转动. (a) 写电场和磁场; (b) 当 $\dot\omega\ne0$ 时求表面感应电场与力矩; (c) 加外力矩和机械转动惯量 $I$, 求等效附加惯量; (d) 常角速度时求磁能内外分配; (e) 求场角动量并核对惯量关系.
\end{exercise}
\begin{solution}
第16章期末题1给(b)、(d)、(e)的详细积分及 $I_\mathrm{mag}=\mu_0Q^2R/(18\pi)$. 此题(a)的静电场为内零、外 $Q\hat{\rr}/(4\pi\eps r^2)$; 磁场为内 $B_0\hat{\vct z}$、外 $B_0R^3(2\cos\theta\hat{\rr}+\sin\theta\hat{\vct\theta})/(2r^3)$, $B_0=\mu_0Q\omega/(6\pi R)$. (c) 感应力矩为 $-I_\mathrm{mag}\dot\omega\hat{\vct z}$, 故 $I\dot\omega=\tau_\mathrm{ext}-I_\mathrm{mag}\dot\omega$, 即 $(I+I_\mathrm{mag})\dot\omega=\tau_\mathrm{ext}$. 所用随时间磁场取准静态近似, 未把辐射反作用包含在等效惯量中.
\end{solution}
\begin{exercise}
原题7 (15分加分). 磁单极模型满足 $\nabla\cdot\bb=\mu_0\rho_m$, $\nabla\times\ee=-\mu_0\jj_m-\partial_t\bb$, 其余两方程为普通电荷方程. 磁荷受力 $q_m(\bb-\vct v\times\ee/c^2)$. 推广坡印廷定理; 若受力叉积项换成加号, 单凭能量守恒能否排除?
\end{exercise}
\begin{solution}
点磁荷的功率是 $\vct F\cdot\vct v=q_m\bb\cdot\vct v$, 叉积项不做功. 单位体积机械功率为 $\jj_e\cdot\ee+\jj_m\cdot\bb$. 把新法拉第定律放入能量推导, 得
$\partial_tu+\nabla\cdot\vct S=-\jj_e\cdot\ee-\jj_m\cdot\bb$, $u,\vct S$ 不变. 若叉积项误换加号, 它仍不做功, 能量定理照样成立, 因而能量守恒不能单独确定该符号; 需要动量或对偶一致性. 磁荷是本题设定的理论模型, 不表示实验已发现磁单极.
\end{solution}
原题1、2、4、5、6只列 Griffiths 8.4、8.6、9.10、9.12、9.13, 不补商业题面.
